\documentclass[10pt,a4paper]{article}
\usepackage[utf8]{inputenc}
\usepackage[T1]{fontenc}
\usepackage[margin=0.62in]{geometry}
\usepackage{times}
\usepackage{microtype}
\usepackage{xcolor}
\PassOptionsToPackage{hyphens}{url}
\usepackage{hyperref}

\hypersetup{
    colorlinks=true,
    linkcolor=blue,
    urlcolor=blue,
    citecolor=blue
}

\begin{document}
\thispagestyle{empty}

% --- Header / Institutional Contact ---
\noindent
\textbf{College of Information Science and Engineering, Huaqiao University}\\
{\small No. 668 Jimei Avenue, Xiamen, Fujian 361021, P.~R.~China \hfill Corresponding Author: Fang Duan, E-mail: \texttt{nkfetsh@gmail.com}}
\vspace{0.4ex}
\hrule height 0.6pt
\vspace{1ex}

\noindent
\textbf{Date:} October 26, 2026 \hfill \textbf{Topic / Section:} Brain--Computer Interfaces \& Neural Signal Processing\\[0.8ex]
\textbf{To:} Prof.~Warren M.~Grill, Editor-in-Chief, \textit{Journal of Neural Engineering} (JNE)\\
\phantom{\textbf{To:}} Department of Biomedical Engineering, Duke University, Durham, NC 27708, USA\\[0.8ex]
\textbf{Subject:} Submission of Research Paper: \textbf{``Physiology-Guided Temporal Fusion for Cross-Subject EEG--fNIRS Brain--Computer Interfaces''}

\vspace{1ex}
\noindent Dear Prof.~Grill and Members of the Editorial Board,

\smallskip
We are pleased to submit our original research manuscript entitled \textbf{``Physiology-Guided Temporal Fusion for Cross-Subject EEG--fNIRS Brain--Computer Interfaces''} for consideration for publication in the \textit{Journal of Neural Engineering} (JNE).

\smallskip
\noindent\textbf{Context \& Neural Engineering Challenge:}
Hybrid EEG--fNIRS interfaces offer a powerful non-invasive window into neural function by pairing high-temporal-resolution electrophysiology with localized capillary hemodynamics. However, developing zero-calibration, cross-subject brain--computer interfaces (BCIs) remains constrained by the fundamental biological timescale disparity between electrophysiological transients (milliseconds) and capillary hemodynamic responses (seconds). When unconstrained deep cross-attention networks are trained in cross-subject leave-one-subject-out (LOSO) regimes, they frequently learn spurious, acausal alignments where sluggish hemodynamic tokens mistakenly attend to future neural activations ($\delta_{ij} < 0$), impairing generalization to unseen participants.

\smallskip
\noindent\textbf{Key Contributions \& Empirical Rigor:}
To resolve this biophysical bottleneck, we propose \textbf{NVC-DelayFormer}, an architecture that embeds biophysical neurovascular timing priors directly into multimodal cross-attention. Our contributions are substantiated through rigorous participant-isolated evaluations and open-science benchmarking:
\begin{enumerate}\setlength\itemsep{0.08em}\setlength\parskip{0pt}
  \item \textbf{Physiology-Guided Cross-Attention Formulation:} Multi-head cross-attention is parameterized by multicenter Gaussian delay priors grounded in canonical hemodynamic response delay windows ($\tau_h \in \{3.5, 4.5, 5.5, 6.5\}$\,s) and token timestamp masks that explicitly eliminate acausal backwards-in-time connections, ensuring temporal precedence across modalities.
  \item \textbf{Factorial Ablation \& Dual-Cohort LOSO Validation:} Under strict leave-one-subject-out (LOSO) protocols across both the primary 29-participant cohort (TU Berlin) and an independent 15-participant external cohort (Buccino et al., 81.03\% accuracy; Supplementary Section~S.V), factorial ablations explicitly confirm cross-hardware generalizability and isolate the distinct marginal contributions of the delay prior versus the causal mask.
  \item \textbf{Unsupervised Recursive Covariance Alignment (CrEA):} To tackle inter-subject covariance drift during test time, we introduce Causal Recursive Euclidean Alignment (CrEA), which recursively updates reference covariances online without session lookahead, reducing relative covariance error from 45.15\% to 29.38\% within the first 8 trials. Across 29 participants and three random seeds, NVC-DelayFormer achieves a grand mean accuracy of \textbf{78.22\% $\pm$ 1.00\%} (primary seed: \textbf{79.37\% $\pm$ 7.73\%}), significantly exceeding standard 2D-EEGNet (72.89\%, $p = 1.10\times 10^{-4}$) and unconstrained cross-attention ($p = 7.18\times 10^{-4}$). In task-specific training, accuracy reaches \textbf{73.24\%} for working memory and \textbf{66.32\%} for motor imagery.
  \item \textbf{Compact Embedded Footprint:} With only \textbf{63,932 trainable parameters} and \textbf{24.79\,M MACs}, the model requires $17.5\times$ fewer parameters than EEG Conformer and operates with a batch-one forward latency of \textbf{2.22\,ms} on an NVIDIA RTX 4060 Laptop GPU, satisfying the operational envelope of wearable neural interfaces.
\end{enumerate}

\smallskip
\noindent\textbf{Mandatory Declarations \& Compliance:}
\begin{itemize}\setlength\itemsep{0.06em}\setlength\parskip{0pt}
  \item \textbf{Originality:} This manuscript represents original work not published previously and not currently under consideration elsewhere.
  \item \textbf{Authorship Approval:} All authors have read and approved the final manuscript and agree with its submission to the \textit{Journal of Neural Engineering}.
  \item \textbf{Ethical Approval:} All human participant analyses used de-identified open-access data approved by the Ethics Committee of the Institute of Psychology and Ergonomics, Technische Universit{\"a}t Berlin (protocol SH\_01\_20150330), in accordance with the Declaration of Helsinki.
  \item \textbf{Conflict of Interest \& Funding:} The authors declare no competing financial or non-financial interests. No specific external grant supported this research.
  \item \textbf{Reproducibility \& Open Science:} Full source code, pre-trained checkpoints, test suites, and audit logs are openly accessible at \href{https://github.com/ujujuj3061831150/nvc-delayformer-bci}{GitHub Repository (ujujuj3061831150/nvc-delayformer-bci)}.
\end{itemize}

\smallskip
\noindent Thank you very much for your time and consideration of our manuscript.

\smallskip
\noindent Sincerely,\\[0.8ex]
\textbf{Fang Duan} (Corresponding Author, E-mail: \texttt{nkfetsh@gmail.com}) \\
\textbf{Hongjie Chen} (First Author, E-mail: \texttt{chenhj976@gmail.com}) \\
College of Information Science and Engineering, Huaqiao University, Xiamen 361021, Fujian, China
\end{document}
